\begin{recetaConImagen}{Guisillo de Carne}{4 personas}{40 min activos · 2 h total}{receta:guisillo_carne}
    {imagenes/guisillo.png}{Cada vez me sale mejor el guisillo}

    \origenReceta{\faHome\ De casa · Receta de familia}

    \begin{minipage}[t]{0.35\textwidth}
        \section*{Ingredientes}
        \begin{itemize}[leftmargin=*]
            \item 600 g de carne (pollo y costilla de cerdo)
            \item 800 g de patatas
            \item 2 zanahorias
            \item 1 cebolla y 2 dientes de ajo
            \item 2 tomates maduros
            \item 2 pimientos (rojo y verde)
            \item 200 ml de vino blanco
            \item 1 hoja de laurel
            \item colorante o azafrán
            \item sal
            \item 1,5 l de caldo de pollo, aprox.
        \end{itemize}
    \end{minipage}
    \hfill
    \begin{minipage}[t]{0.6\textwidth}
        \section*{Preparación}
        \begin{enumerate}
            \item \textbf{Sellar la carne:} En una olla con aceite, dorar los trozos de pollo y costilla.
            \item \textbf{Sofrito:} Añadir por orden: 1º el ajo y la cebolla, 2º los pimientos, 3º la zanahoria y 4º el tomate.
            \item \textbf{Reducción:} Cuando no haya caldo del sofrito, verter un vaso de vino blanco y esperar a que reduzca.
            \item \textbf{Cocción:} Añadir caldo de pollo y/o agua hasta cubrir, el laurel, el colorante y la sal cuando hierva. Cocer a fuego suave durante 60-90 minutos, hasta que la carne esté tierna.
            \item \textbf{Finalizar:} Añadir las patatas y cocer otros 20 minutos, hasta que estén tiernas.
        \end{enumerate}
    \end{minipage}

    \vspace{1em}
    \begin{tcolorbox}[colback=orange!5, colframe=orange!50, title=Sugerencia, size=small]
        Corta las patatas "chasqueándolas" para que el caldo espese de forma natural.
    \end{tcolorbox}
\end{recetaConImagen}
